\documentclass[11pt]{article}

\usepackage[a4paper,margin=2.4cm]{geometry}
\usepackage[T1]{fontenc}
\usepackage[utf8]{inputenc}
\usepackage{lmodern}
\usepackage{amsmath}
\usepackage{amssymb}
\usepackage{booktabs}
\usepackage{array}
\usepackage{longtable}
\usepackage{enumitem}
\usepackage{xcolor}
\usepackage{graphicx}
\usepackage[hidelinks]{hyperref}
\usepackage{microtype}
\emergencystretch=3em

\definecolor{codegray}{gray}{0.94}
\newcommand{\code}[1]{{\setlength{\fboxsep}{1.5pt}\colorbox{codegray}{\small\texttt{#1}}}}
\newcommand{\file}[1]{\texttt{#1}}

\setlist{itemsep=2pt, topsep=4pt}
\setcounter{tocdepth}{2}

\title{\textbf{nuScenes $\rightarrow$ OpenDRIVE $\rightarrow$ CARLA / RoadRunner}\\[6pt]
\large Comprehensive Progress Report:\\ Rebuilding Real Driving Scenes in Simulation}
\author{Keshu W.}
\date{July 22, 2026}

\begin{document}
\maketitle

\begin{abstract}
\noindent
This report documents, end to end, a pipeline that reconstructs real recorded driving
scenes from the \textbf{nuScenes v1.0-mini} dataset inside two simulators. All ten mini
scenes --- spanning the maps \emph{boston-seaport}, \emph{singapore-onenorth},
\emph{singapore-queenstown}, and \emph{singapore-hollandvillage} --- are converted into
OpenDRIVE road networks with real lane markings and traffic lights, enriched with ground,
OpenStreetMap buildings, and assembled signal hardware, and then replayed --- ego vehicle
plus all 703 recorded traffic participants --- in \textbf{CARLA 0.9.15} and in
\textbf{MathWorks RoadRunner R2026a / RoadRunner Scenario}. Validation is systematic:
per-scene real-vs-simulated six-camera montage GIFs, real-vs-simulated LiDAR bird's-eye
views, reconstructed top-down renders, and OSM-georeferenced overlay figures. As of this
writing, every scene has a converted and RoadRunner-hardened \file{.xodr} map, a full
set of CARLA renders, a built RoadRunner ``city'' scene, and a saved collision-safe
replay scenario. The report describes the architecture, every processing stage and its
algorithms, the data schemas, the produced artifacts with per-scene statistics, and the
known limitations.
\end{abstract}

\tableofcontents
\clearpage

%=====================================================================
\section{Introduction}

\subsection{Objective}
The objective is to take a real, annotated driving log --- a nuScenes \emph{scene},
i.e.\ roughly 20 seconds of urban driving with an HD vector map and 2\,Hz cuboid
annotations for every traffic participant --- and reproduce it in simulation: the same
road network with its true lane markings and traffic-light placement, comparable
surroundings (ground plane, buildings), and the same traffic participants moving along
their recorded trajectories around a controllable or replayed ego vehicle. Two simulator
back-ends are fed from one simulator-neutral intermediate representation, enabling both
game-engine experimentation (CARLA) and scenario-editing workflows (RoadRunner
Scenario).

\subsection{Dataset and scope}
The source data is \textbf{nuScenes v1.0-mini}: 10 scenes, 404 keyframes total
($\sim$39--41 keyframes per scene at 2\,Hz), recorded in Boston and Singapore across
four map locations. Each map ships an HD vector map (lanes, lane connectors, dividers,
stop lines, pedestrian crossings, traffic lights) with a published WGS84 origin, which
this project exploits for georeferencing throughout.

\subsection{Pipeline at a glance}
\begin{center}
\scriptsize
\setlength{\tabcolsep}{4pt}
\begin{tabular}{@{}llll@{}}
\toprule
\textbf{Stage} & \textbf{Tool} & \textbf{Env} & \textbf{Output} \\
\midrule
1. Map conversion & \file{convert\_nuscenes\_to\_xodr.py} & Py 3.12 (\file{nusc2xodr}) & \file{.xodr} + \file{\_meta.json} \\
2. RoadRunner hardening & \file{fix\_xodr\_for\_roadrunner.py} & Py 3.12 & repaired \file{.xodr} \\
3. CARLA load \& replay & \file{load\_xodr\_in\_carla.py} & Py 3.7 (\file{carla915}) & live sim, cam/LiDAR renders \\
4. Validation & 2 notebooks + 5 scripts & Py 3.12 & montages, GIFs, BEV, overlays \\
5. RoadRunner city scene & \file{nuscenes\_ground\_buildings.m} & MATLAB + Scene Builder & \file{*\_city.rrscene} \\
6. RoadRunner replay & \file{replay\_agents.m} & MATLAB + RR Scenario & \file{*\_replay.rrscenario} \\
\bottomrule
\end{tabular}
\end{center}

The intermediate representation is an OpenDRIVE 1.4 file plus a JSON \emph{sidecar}
(\file{\_meta.json}) carrying everything OpenDRIVE cannot: the local-origin offset, the
patch bounds, the ego trajectory, and the per-frame poses of every annotated agent
(Section~\ref{sec:sidecar}).

\subsection{Repository layout}
\begin{itemize}
\item \file{nuscenes2xodr/} --- the Python toolchain (converter, fixer, CARLA loader,
      notebooks, plot/composite scripts, two conda environment files) and its
      \file{output/} tree with all \file{.xodr} maps and rendered artifacts.
\item \file{output\_roadrunner/} --- RoadRunner-hardened \file{.xodr} + sidecars, one
      \file{.osm} extract per scene, and a per-scene subfolder with the generated HD map
      (\file{\_city.rrhd}), georeferenced \file{.xodr} copy, marking cache, and agent
      trajectory CSVs.
\item \file{nuscenes\_test1/} --- the RoadRunner project (Assets, Scenes, Scenarios).
\item \file{v1.0-mini/} --- a local copy of the nuScenes mini split (49k+ files).
\item Five MATLAB scripts at the root: \file{insert\_intersection.m},
      \file{load\_nuscenes\_xodr.m}, \file{load\_osm\_roads.m},
      \file{nuscenes\_ground\_buildings.m}, \file{replay\_agents.m}.
\end{itemize}

\subsection{Execution environments}
Two conda environments are required on the Python side because CARLA 0.9.15 ships only a
Python 3.7 binding (\file{carla-0.9.15-cp37-cp37m}); there is no \file{carla} wheel for
Python 3.12:
\begin{itemize}
\item \file{nusc2xodr} (Python 3.12): \file{numpy}, \file{shapely},
      \file{nuscenes-devkit}. Runs the converter, fixer, notebooks, and plot scripts.
\item \file{carla915} (Python 3.7): \file{numpy}, \file{pygame}, plus the bundled CARLA
      wheel installed manually. Runs the loader against a running
      \file{CarlaUE4.exe} server.
\end{itemize}
The MATLAB side requires Automated Driving Toolbox (for \code{roadrunner} /
\code{roadrunnerHDMap}), a RoadRunner license with \textbf{Scene Builder} (HD-map
import auto-builds junctions/props), and a \textbf{RoadRunner Scenario} license for the
replay scenarios. RoadRunner is launched programmatically from MATLAB
(\file{C:\textbackslash Program Files\textbackslash RoadRunner R2026a}), with session
reuse: every script probes an existing \code{rrApp} handle before launching a new
instance, and stops any active simulation before changing scenes (dataset changes fail
while RoadRunner is simulating).

%=====================================================================
\section{Coordinate Frames and Georeferencing}
\label{sec:frames}

\textbf{Map frame.} Raw nuScenes coordinates (metres, right-handed, $x$ east / $y$
north per map).

\textbf{Patch-local frame.} Everything exported is shifted by
$\mathit{origin} = (x_{\min}, y_{\min})$, the lower-left corner of the scene's region
patch, keeping geometry near the origin for numerical stability. No rotation, no axis
flip.

\textbf{WGS84.} Each map's published origin (latitude/longitude of map coordinate
$(0,0)$) is used with an equirectangular linearization
($R_\oplus = 6\,378\,137$\,m):
\[
\mathrm{lat} = \mathrm{lat}_0 + \frac{y}{R_\oplus}\cdot\frac{180}{\pi},
\qquad
\mathrm{lon} = \mathrm{lon}_0 + \frac{x}{R_\oplus \cos(\mathrm{lat}_0)}\cdot\frac{180}{\pi}.
\]
The published origins used consistently across Python and MATLAB:

\begin{center}
\small
\begin{tabular}{lrr}
\toprule
\textbf{Map} & \textbf{lat$_0$} & \textbf{lon$_0$} \\
\midrule
boston-seaport            & 42.336849169438615 & $-$71.05785369873047 \\
singapore-onenorth        & 1.2882100868743724 & 103.78475189208984 \\
singapore-hollandvillage  & 1.2993652317780957 & 103.78217697143555 \\
singapore-queenstown      & 1.2782562240223188 & 103.76741409301758 \\
\bottomrule
\end{tabular}
\end{center}

\textbf{CARLA convention.} CARLA is left-handed and flips $Y$ when ingesting OpenDRIVE;
the loader converts every pose with $(x, y, \psi) \rightarrow (x, -y, -\psi)$.

\textbf{RoadRunner convention.} Both the OpenDRIVE and the HD-map imports are placed by
\emph{projection}: the files carry a real transverse-mercator \code{geoReference}
(a \code{+proj=tmerc} string) centred on the patch corner, the
scene world origin is set to the same point via \code{changeWorldSettings}, and imports
use \code{ProjectionMode="FullProjection"} --- so multiple layers land in registration
deterministically instead of being centred on their data.

%=====================================================================
\section{Stage 1 --- nuScenes to OpenDRIVE Converter}
\label{sec:converter}

\file{convert\_nuscenes\_to\_xodr.py} converts the map region covered by one scene into
OpenDRIVE 1.4. It runs standalone on \file{nuscenes-devkit}/\file{shapely}/\file{numpy}.

\subsection{Core algorithm}
\begin{enumerate}
\item \textbf{Ego trajectory and patch.} Walk the scene's sample chain; for each
      keyframe read the \code{LIDAR\_TOP} sample's \code{ego\_pose}, yielding
      $(x, y, \psi)$ poses (yaw extracted from the quaternion). The region patch is the
      bounding box of the trajectory padded by \code{--padding} (default 50\,m); an
      explicit \code{--region} override is supported.
\item \textbf{Lane selection.} \code{get\_records\_in\_patch} selects all \code{lane}
      and \code{lane\_connector} records intersecting the patch.
\item \textbf{Geometry and width.} Each centerline is discretized from its arcline path
      at \code{--resolution} (default 0.5\,m). Per-lane width is estimated as
      $\text{polygon area} / \text{centerline length}$, clamped to $[2, 6]$\,m
      (fallback 3.5\,m), unless \code{--lane-width} fixes it globally.
\item \textbf{Topology.} A union-find clusters lane connectors that share an incoming
      \emph{or} outgoing lane; each cluster becomes one \code{<junction>}. Because
      nuScenes connectors are strictly one-in/one-out, every lane has at most one
      junction per end, giving well-defined predecessor/successor links.
\item \textbf{Emission.} Each lane and connector becomes a \code{<road>}; connectors
      carry the junction id. The plan view is a chain of straight \code{<line>}
      geometries (no arcs/spirals). Each road has a single driving lane with id $-1$;
      a \code{laneOffset} of $+w/2$ shifts the lane reference so the driving lane
      \emph{straddles} the true centerline. Junctions emit one \code{<connection>} per
      connector with \code{<laneLink from="-1" to="-1">}.
\end{enumerate}

\subsection{Lane markings}
Marking types come from the real nuScenes divider data: for each lane, the dominant
non-NIL \code{segment\_type} among its left/right divider segments is mapped to an
OpenDRIVE \code{roadMark} --- \code{DASHED} $\rightarrow$ broken, \code{SOLID} (and,
approximately, \code{ZIGZAG}) $\rightarrow$ solid, a \code{DOUBLE} qualifier doubles the
type, and a \code{YELLOW} substring selects yellow. The mark on the centerline side is
carried by the center lane, the outer mark by the driving lane (width 0.13\,m,
standard weight). Lane connectors emit no markings, matching unmarked junction
interiors in the source map.

\subsection{Optional layers}
\begin{itemize}
\item \textbf{Crosswalks} (\code{--crosswalks}): \code{ped\_crossing} polygons are
      exported both into the sidecar and as OpenDRIVE
      \code{<object type="crosswalk">} outlines attached to the nearest road via
      projection to $(s,t)$ coordinates.
\item \textbf{Buildings} (\code{--buildings}): nuScenes has no building layer, so
      footprints are \emph{synthesized} in the negative space: the union of occupied
      layers (drivable area, road segments, walkways, carparks, crossings, lanes,
      buffered 2\,m) is subtracted from the patch; leftover blocks (area
      $\ge 120$\,m$^2$, 3\,m setback) are tiled on a 14\,m grid with
      $11.2\times11.2$\,m boxes and deterministic pseudo-random heights of 8--24\,m.
      These go to the sidecar only.
\item \textbf{Traffic lights} (\code{--traffic-lights}): OpenDRIVE \code{<signal>}
      elements (type 1000001, dynamic, $z$-offset 5\,m) are placed at the end of each
      approach lane and grouped into per-phase \code{<controller>}s by approach axis
      --- CARLA's importer turns these into functional, cycling
      \code{carla.TrafficLight} groups. With \code{--osm-signals} the converter fetches
      the real OSM \code{traffic\_signals} nodes for the region from the Overpass API
      (two mirrors, retries) and places one signal per real node on the nearest
      approach lane, so the light \emph{count} matches reality; without it, every
      junction is lit.
\end{itemize}

\subsection{The metadata sidecar}
\label{sec:sidecar}
Each map ships with \file{<name>\_meta.json}:

\begin{center}
\small
\begin{tabular}{@{}lp{9.6cm}@{}}
\toprule
\textbf{Key} & \textbf{Content} \\
\midrule
\code{scene}, \code{map} & scene name and map location \\
\code{origin} & $\{x, y\}$ patch corner in the map frame (the local-frame origin) \\
\code{patch} & $\{x_{\min}, y_{\min}, x_{\max}, y_{\max}\}$ region bounds \\
\code{ego\_start\_local} & first ego pose $\{x, y, \mathit{yaw}\}$, local frame \\
\code{ego\_trajectory\_local} & list of ego poses, one per keyframe \\
\code{frame\_dt} & keyframe interval, 0.5\,s \\
\code{agent\_frames} & per keyframe: list of agents $\{$\code{id} (instance token),
  \code{category}, \code{x}, \code{y}, \code{yaw}, \code{wlh}$\}$; dynamic categories
  (\code{vehicle.*}, \code{human.pedestrian.*}) by default,
  \code{--include-static} adds cones/barriers/debris \\
\code{buildings} & synthesized footprints $\{x, y, s_x, s_y, \mathit{yaw}, \mathit{height}\}$ \\
\code{crosswalks} & local-frame polygons \\
\code{counts} & $\{$lanes, connectors, junctions$\}$ \\
\bottomrule
\end{tabular}
\end{center}

The ego trajectory and \code{agent\_frames} are aligned one-to-one, so a frame index
identifies the same instant in every downstream artifact (camera frames, LiDAR sweeps,
BEV renders, RoadRunner trajectories).

%=====================================================================
\section{Stage 2 --- Hardening OpenDRIVE for RoadRunner}
\label{sec:fixer}

The raw converter output loads in CARLA but is rejected or mishandled by RoadRunner.
\file{fix\_xodr\_for\_roadrunner.py} batch-processes \file{nuscenes2xodr/output/*.xodr}
into \file{output\_roadrunner/} (copying the sidecars along) and applies six fixes:

\begin{enumerate}
\item \textbf{Endpoint snapping.} Connector endpoints within 2\,m of their linked
      lane's endpoint are snapped exactly onto it, closing sub-2\,m gaps that leave
      roads disconnected in RoadRunner.
\item \textbf{Plan-view simplification.} Each road polyline is reduced with
      Douglas--Peucker at 0.03\,m lateral tolerance and the plan view re-emitted ---
      the 0.5\,m discretization otherwise creates hundreds of control points per road
      that RoadRunner struggles to manage. (This is why the hardened files are
      3--7$\times$ smaller, e.g.\ scene-0103: 912\,kB $\rightarrow$ 199\,kB.)
\item \textbf{Marking defaults.} Real lanes whose divider data was NIL get standard
      marks (broken white centerline, solid white edge); junction interiors stay
      unmarked.
\item \textbf{Lane-level links.} Connector roads receive lane-level
      \code{<predecessor>}/\code{<successor>} links mirroring the road-level links ---
      RoadRunner derives maneuver topology from lane links.
\item \textbf{Real header bounds.} North/south/east/west are recomputed from actual
      geometry ($\pm$10\,m pad); the converter's \code{south=west=0} caused
      RoadRunner's ``Bounds of file invalid'' rejection.
\item \textbf{Real geoReference.} The header's projection string is rewritten to a
      transverse-mercator centred on the scene origin's true latitude/longitude
      (Section~\ref{sec:frames}), so GIS layers and aerial imagery line up.
\end{enumerate}

%=====================================================================
\section{Stage 3 --- CARLA Loading and Replay}
\label{sec:carla}

\file{load\_xodr\_in\_carla.py} ($\sim$1{,}470 lines) targets a running CARLA 0.9.15
server. The world is generated with
\code{client.generate\_opendrive\_world()} (vertex distance 2\,m, max road length
500\,m, extra lane width 0.6\,m, smoothed junctions). The recorded ego start pose is
converted through the $Y$-flip and \emph{snapped} onto the nearest drivable lane via
\code{map.get\_waypoint(project\_to\_road=True)}; the ego (red Tesla Model~3) spawns
there. Scenarios are selectable by id (0--9 in nuScenes order), scene number, or
explicit path; \code{--all-scenarios} runs the whole mini set back-to-back with
per-scenario error isolation (used for batch camera/LiDAR rendering). A weather preset
is applied (\code{ClearNoon} default).

\subsection{Driving modes}
\begin{description}
\item[Replay (default).] Deterministic synchronous-mode playback. The fixed step is
      the keyframe interval over the substep count (default 0.05\,s); the ego has physics
      disabled and is teleported each substep along a Catmull--Rom interpolation of the
      keyframes (shortest-arc yaw interpolation), giving even velocity and no camera
      flicker. Ground height is queried through a 1\,m-grid cache for even frame
      timing, and every actor is grounded by its own bounding box (the offset
      $e_z - b_z$ places the box bottom at road level --- pedestrians don't sink,
      vehicles don't float). The chase camera's yaw is exponentially smoothed.
      Playback loops by default; asynchronous mode and ego physics are restored on
      exit.
\item[Follow.] The ego is driven \emph{physically} by CARLA's
      \code{VehiclePIDController} (lateral $K_P{=}1.0$, $K_I{=}0.1$; longitudinal
      $K_P{=}1.0$, $K_I{=}0.05$; $dt=0.05$\,s) over the recorded waypoints, with
      per-segment target speed $\min(60, \|\Delta p\|/\mathit{dt}\cdot3.6)$\,km/h, a
      3\,m arrival radius, and a 6\,s per-target budget to escape stalls. The recorded
      agents are kept \emph{in step with the ego's actual progress}: the ego's
      projection onto the current route segment yields a fractional keyframe time that
      drives the agent replayer. A red GOAL marker (vertical line + label) shows the
      current target.
\item[Manual.] A $1280\times720$ pygame window with chase camera (FOV 90) and HUD
      (speed, gear, steer, autopilot state, agent/frame counters). Keyboard (W/S/A/D,
      Space handbrake, Q reverse, P autopilot toggle, ESC quit) or gamepad/wheel
      (\code{--joystick}: axis-mapped steer/throttle/brake with deadzone). The recorded
      agents replay around the driver in real time, looped.
\item[Autopilot.] Traffic Manager drives the ego on the converted map (speed offset
      30\%, ignoring the scripted lights/signs).
\end{description}

\subsection{Agent replay engine}
A shared \code{\_AgentReplayer} spawns each annotated agent on first appearance and
teleports it (physics off) with linear position and shortest-arc yaw interpolation
between keyframes; agents absent from a keyframe are hidden 200\,m underground rather
than destroyed. nuScenes categories map to CARLA blueprints by kind (pedestrian walkers;
bicycles: Crossbike/Century/Omafiets; motorcycles: Low Rider/YZF/Ninja/Vespa; bus:
Fuso Rosa; truck: CarlaCola/Firetruck; cars: A2/TT/Model3/Patrol/Prius/Leon/C3).
Blueprint and color choice are keyed on a stable hash of the nuScenes instance token, so
an agent keeps its model and color across frames \emph{and} across runs; red is
reserved for the ego. Spawn retries lift the actor $+1$/$+3$\,m on collision.
\code{--max-agents} caps the population by order of first appearance;
\code{--replay-speed} scales time.

\subsection{Trajectory overlay}
During replay/follow, past and future trajectories are drawn for the ego and every
agent with track data, sampled at \code{--traj-hz} (default 2\,Hz): 2\,s of history
(dimmed to 35\%) and 6\,s of future (full brightness), as short-lived debug lines and
points re-drawn each keyframe. Colors: ego red, vehicles blue, pedestrians green,
cyclists/motorcyclists magenta.

\subsection{Sensors: six-camera rig and LiDAR}
\textbf{Calibrated rig (default).} When a scene's calibration sidecar exists (produced
by a companion experiment, \file{phase0\_calib\_gt}), the six nuScenes cameras are
spawned \emph{free-floating} --- CARLA's vehicle pivot differs from the nuScenes ego
origin, so attachment would bias extrinsics. Each tick every camera is teleported to
$\mathit{ego\_pose}\circ\mathit{extrinsic}$ (rotation composed in matrix form, then
converted back to CARLA Euler angles) and renders at the native $1600\times900$ with
per-camera FOV $=2\arctan\!\big(W/2f_x\big)$ from the real intrinsics (e.g.\
CAM\_FRONT $64.56^\circ$ rather than the approximate $70^\circ$). Saved views therefore
share the real cameras' projection and are directly usable for calibration-sensitive
perception evaluation (BEVFormer-class models).

\textbf{Approximate rig (fallback / \code{--approx-rig}).} Six cameras attached at
roughly the nuScenes mounts (front trio at $z=1.51$\,m, yaws $0/\pm55^\circ$, FOV 70;
back at $180^\circ$, FOV 110; back sides at $\pm110^\circ$, FOV 70), default
$800\times450$.

Frames are saved once per keyframe into a per-scene, per-camera folder tree under
\file{output/cameras/} (follow mode writes to \file{output/cameras\_follow/} so runs
don't clash), mirroring the nuScenes per-camera layout.

\textbf{LiDAR.} \code{--lidar} attaches an HDL-32E-like \code{sensor.lidar.ray\_cast}
at the nuScenes mount ($x=0.90$, $z=1.84$\,m; $+10^\circ/-30^\circ$ vertical FOV, 32
channels, 70\,m range, $\sim$34k points per sweep). In synchronous mode the rotation
rate is locked to the tick so each keyframe stores exactly one full $360^\circ$ sweep as
an ASCII \file{.ply} under \file{output/lidar\_carla/<scene>/}.

\subsection{Scenery drawing}
Sidecar buildings can be realized as \emph{solid} blocks by tiling a static-prop box
mesh (shipping container by default; prop size probed at runtime) up to
\code{--building-layers} high, or as persistent wireframe boxes. Sidecar crosswalk
polygons are drawn as white zebra stripes (0.5\,m bars, 0.6\,m gaps) slightly above the
road, since CARLA's standalone OpenDRIVE mode paints no crossings.

%=====================================================================
\section{Stage 4 --- Validation and Visualization}
\label{sec:validation}

\subsection{Real vs.\ simulated cameras}
\file{align\_sim\_vs\_real.ipynb} builds zero-padding, labelled $2\times3$ six-camera
montages (front-left/front/front-right over back-left/back/back-right, $480\times270$
cells) for each keyframe --- real images from the devkit on top, CARLA renders below ---
and animates three GIFs per scene at the native 2\,Hz: \file{nuscenes\_real.gif},
\file{carla\_sim.gif}, and \file{comparison.gif} (REAL banner over SIM banner). A
\code{MODE} switch selects the replay or follow camera tree; missing simulated cameras
render as placeholders so the real side always works.

\subsection{OSM georeferencing check}
\file{align\_osm.ipynb} projects each scene's patch and ego trajectory to WGS84
(Section~\ref{sec:frames}), fetches OSM roads, buildings, and \code{traffic\_signals}
for the bounding box from Overpass (stdlib-only HTTP, three mirror endpoints), and plots
OSM (grey) under the nuScenes lane centerlines (blue), ego trajectory (red), and real
signal nodes (yellow). The resulting figures
(\file{output/osm\_aligned/<scene>.png}) validate both the georeferencing used
everywhere and the \code{--osm-signals} light counts.

\subsection{Bird's-eye-view and LiDAR renderers}
\begin{itemize}
\item \file{plot\_topdown.py} re-renders the scene \emph{from the exported artifacts
      alone} (\file{.xodr} lane surfaces + sidecar poses): ego-centred, north-up BEV
      frames (100\,m window) with oriented agent boxes, heading ticks, and the same
      history/future trajectory color scheme as CARLA; per-scene GIF at 4\,fps.
\item \file{plot\_lidar.py} renders the original nuScenes \code{LIDAR\_TOP} sweeps in
      the same ego-centred north-up frame (height-colored, $-2$ to $8$\,m clip, dark
      background).
\item \file{plot\_lidar\_carla.py} renders the CARLA-captured sweeps identically
      (left$\rightarrow$right-handed flip, $+1.84$\,m mount height, rotation by the
      recorded ego yaw), enabling a direct real-vs-sim LiDAR comparison.
\item The composite builders \file{build\_aligned\_topdown.py} and
      \file{build\_aligned\_full.py} assemble per-frame panels: cameras + BEV, and the
      full three-column panel (REAL/SIM cameras, real-over-sim LiDAR, top-down), as
      404 PNG frames and 10 GIFs each.
\end{itemize}

\subsection{Rendered-artifact inventory}
All ten scenes are fully rendered on the CARLA side; the folders below live in the
converter's \file{output/} tree:

\begin{center}
\small
\setlength{\tabcolsep}{4pt}
\begin{tabular}{@{}llr@{}}
\toprule
\textbf{Folder} & \textbf{Content} & \textbf{Files} \\
\midrule
\file{cameras/}          & replay-mode 6-cam frames (10 scenes + 1 backup) & 2{,}658 PNG \\
\file{cameras\_follow/}  & follow-mode 6-cam frames                        & 2{,}424 PNG \\
\file{aligned/}          & real/sim/comparison GIFs (replay)               & 30 GIF \\
\file{aligned\_follow/}  & real/sim/comparison GIFs (follow)               & 30 GIF \\
\file{lidar/}            & real LiDAR BEV frames + GIFs                    & 404 PNG + 10 GIF \\
\file{lidar\_carla/}     & CARLA LiDAR sweeps, BEV frames, GIFs            & 404 PLY + 404 PNG + 10 GIF \\
\file{topdown/}          & reconstructed BEV frames + GIFs                 & 404 PNG + 10 GIF \\
\file{aligned\_topdown/} & cameras+BEV composites                          & 404 PNG + 10 GIF \\
\file{aligned\_full/}    & cameras+LiDAR+BEV composites                    & 404 PNG + 10 GIF \\
\file{osm\_aligned/}     & OSM overlay figures                             & 10 PNG \\
\bottomrule
\end{tabular}
\end{center}

%=====================================================================
\section{Stage 5 --- RoadRunner City Scenes (MATLAB)}
\label{sec:rrscene}

The RoadRunner side was developed in four increments.

\subsection{Programmatic building blocks}
\file{insert\_intersection.m} established the workflow: since the MATLAB API cannot edit
scene geometry directly, two crossing two-lane roads are described as a
\code{roadrunnerHDMap} (lane boundaries + lanes with travel directions); on import,
Scene Builder's \emph{overlap groups} auto-construct the four-way junction (surface,
corner radii, markings). \file{load\_nuscenes\_xodr.m} then established direct OpenDRIVE
import of a converted scene.

\subsection{Deterministic multi-layer registration}
\file{load\_osm\_roads.m} co-registers the nuScenes road network with the OSM road
network of the same area. The OSM extract is converted to OpenDRIVE through MATLAB's
\code{drivingScenario} importer/exporter; both OpenDRIVE files then receive a real
transverse-mercator \code{geoReference} at their local origins (injected textually into
the header); the RoadRunner world origin is set to the nuScenes patch corner; and both
files are imported with \code{FullProjection}. A diagnostic overlay figure predicts the
in-scene registration and is used to calibrate a per-scene \code{manualShift}
(east/north metres) absorbing OSM's few-metre urban absolute error; the shift is baked
into the OSM file's declared origin so the projection itself carries the correction.

\subsection{The city-scene builder}
\file{nuscenes\_ground\_buildings.m} ($\sim$790 lines) assembles, per scene, a complete
``city'' scene from five sources, all delivered through \emph{one} HD-map
(\file{.rrhd}) import on top of the OpenDRIVE roads:

\begin{enumerate}
\item \textbf{Roads}: the hardened scene \file{.xodr}, imported with signals off
      (signals are assembled separately, below).
\item \textbf{Real markings}: \code{lane\_divider} and \code{road\_divider} polylines
      are extracted from the nuScenes map-expansion JSON (patch-clipped, cached to a
      \file{.mat}), and emitted as HD-map \emph{curve markings} referencing RoadRunner
      marking assets (dashed single white / solid double yellow) drawn 2\,cm above the
      road. They come from the same map frame as the lanes, so they are aligned by
      construction.
\item \textbf{Ground}: a large level slab (patch $+120$\,m margin) built as parallel
      12\,m flat \code{Sidewalk} lanes at $z=-0.15$\,m.
\item \textbf{Buildings}: OSM footprints (the extract is auto-downloaded from the OSM
      API per scene if missing, with a $\sim$200\,m padded bbox; failures are
      non-fatal). Heights come from OSM \code{height} or
      \code{building:levels}$\times$3.2\,m (default 10\,m). Each footprint is
      approximated by oriented boxes: the minimum-area rectangle is used directly when
      it fits well, otherwise the polygon is recursively split across the rectangle's
      long axis (up to depth 3, fill-ratio threshold 0.55) so L/U-shaped buildings
      become several tight boxes. Boxes that cover road are shrunk about their center
      until they clear the lane centerlines by 2\,m (never below 80\% scale ---
      buildings shrink but survive). Boxes are realized as \code{StaticObject}s cycling
      through three downtown building meshes.
\item \textbf{Traffic signals}: positions and facings are taken (configurably) from the
      \file{.xodr} \code{<signal>} entries, the OSM signal nodes, or the surveyed
      nuScenes \code{traffic\_light} layer. Since the recorded point marks the
      \emph{light} (over the road), not the pole, the builder: snaps the facing to the
      nearest lane travelling toward the light; determines the off-road side by probing
      road-sample distances; walks the pole position outward until it clears the
      drivable area; then assembles --- like the stock \file{FourWaySignal} scene --- a
      9.2\,m post at the roadside, a 7.8\,m mast arm spanning back over the road, and
      two three-light heads hung under the arm (outermost over the recorded light
      position), all as \code{StaticObject} boxes with measured FBX half-extents so the
      assets keep natural proportions.
\end{enumerate}

A preview figure (roads, footprints, placed boxes, markings, signal poles/arms/facings)
is rendered before the import for visual QA. Scene selection is generic
(\code{scene-01}\ldots\code{scene-10} or real ids); all ten scenes were built with the
same script into \file{<scene>\_city.rrscene}.

%=====================================================================
\section{Stage 6 --- Agent Replay in RoadRunner Scenario}
\label{sec:rrreplay}

\file{replay\_agents.m} converts the sidecar's \code{agent\_frames} into a saved,
collision-safe RoadRunner Scenario on top of the corresponding city scene.

\subsection{Trajectory export and asset mapping}
Per-agent tracks are collected across frames and written as RoadRunner ``CSV
Trajectory'' files (\code{time,x,y,yaw}; yaw in radians), with
\code{SpawnTime}/\code{RemoveTime} taken from when each agent enters/leaves the
recording. Tracks shorter than 2 frames are skipped. A 2\,mm-per-frame creep along the
heading keeps consecutive waypoints distinct --- the trajectory follower rejects
coincident points, which parked vehicles and stop phases would otherwise produce ---
while preserving stop timing visually. Category mapping: cars cycle through
Sedan/SUV/CompactCar for variety; truck $\rightarrow$ PickupTruck; bus $\rightarrow$
SchoolBus; construction $\rightarrow$ UtilityTruck; emergency $\rightarrow$ Ambulance;
bicycle/motorcycle $\rightarrow$ NCAP two-wheelers; pedestrians (adult/child)
$\rightarrow$ \file{.rrchar} character assets (required: prop/mesh assets abort the
simulator). The ego is a green (\#00B050) sedan.

\subsection{Collision pre-pruning}
RoadRunner Scenario \emph{fails} a simulation on any actor contact, and real annotated
boxes plus substitute assets of different sizes do touch. The script therefore
pre-verifies the entire replay:
\begin{itemize}
\item All accepted tracks (ego first, then agents in order of appearance) are resampled
      on a fine 0.1\,s timeline with \emph{linear} interpolation --- exactly how
      RoadRunner moves actors between waypoints, so contacts \emph{between} the 0.5\,s
      keyframes are caught. Yaw is unwrapped before interpolation.
\item Each candidate is tested against every accepted actor with a separating-axis
      oriented-bounding-box test. Footprints are per-category (e.g.\ car
      $5.1\times2.2$\,m, bus $11.5\times2.9$\,m, pedestrian $0.6\times0.6$\,m) and
      inflated to be strictly more conservative than RoadRunner's own test: vehicle
      half-extents $\times1.25 + 0.3$\,m; pedestrians $+0.15$\,m only (tight boxes so
      groups walking together survive).
\item Actors waiting on their spawn delay are modelled as parked at their first
      waypoint from $t=0$ (they are physically present in RoadRunner).
\item On first contact the agent is \emph{truncated} just before the contact time
      (\code{RemoveTime} takes it out of simulation there); if the contact exists from
      its first frame or the remainder is too short, the agent is dropped. The ego is
      untouchable.
\end{itemize}

\subsection{Scenario assembly}
The script reuses a live RoadRunner session, stops any active simulation, opens the
scene's city scene --- \emph{auto-building it} via
\file{nuscenes\_ground\_buildings.m} if missing --- creates a new scenario, imports the
ego and every surviving agent CSV with its per-actor asset and spawn/remove times,
saves \file{<scene>\_replay.rrscenario}, and echoes any RoadRunner import errors or
warnings. All ten scenarios are saved and playable from the Simulation tool.

%=====================================================================
\section{Results and Current Status}
\label{sec:results}

\subsection{Per-scene statistics}
Statistics computed from the actual artifacts (sidecars, hardened \file{.xodr} files,
and exported trajectories). ``Traj.\ CSV'' is the number of agent trajectories that
survived collision pruning and were imported into the RoadRunner scenario (the ego adds
one more per scene); the difference to ``Agents'' consists of dropped/truncated-away
tracks, single-frame tracks, and unmapped categories.

\begin{center}
\footnotesize
\begin{tabular}{@{}llrrrrrrr@{}}
\toprule
\textbf{Scene} & \textbf{Map} & \textbf{Frames} & \textbf{Agents} & \textbf{Roads} & \textbf{Junc.} & \textbf{Signals} & \textbf{Traj.\ CSV} \\
\midrule
scene-0061 & singapore-onenorth       & 39 &  89 &  73 & 18 & 4 &  57 \\
scene-0103 & boston-seaport           & 40 & 113 &  93 & 18 & 7 & 109 \\
scene-0553 & boston-seaport           & 41 &  40 &  54 &  9 & 9 &  29 \\
scene-0655 & boston-seaport           & 41 & 117 &  77 & 14 & 0 &  73 \\
scene-0757 & boston-seaport           & 41 &  24 &  30 &  5 & 1 &  12 \\
scene-0796 & singapore-queenstown     & 40 &  50 & 130 & 18 & 5 &  35 \\
scene-0916 & singapore-queenstown     & 41 &  94 &  50 & 15 & 0 &  63 \\
scene-1077 & singapore-hollandvillage & 41 &  59 & 137 & 30 & 3 &  38 \\
scene-1094 & singapore-hollandvillage & 40 &  85 &  94 & 18 & 2 &  65 \\
scene-1100 & singapore-hollandvillage & 40 &  32 &  45 &  5 & 8 &  23 \\
\midrule
\textbf{Total} & & 404 & 703 & 783 & 150 & 39 & 504 \\
\bottomrule
\end{tabular}
\end{center}

\subsection{Agent population by category}
Unique annotated agents across all ten scenes:

\begin{center}
\footnotesize
\begin{tabular}{@{}lr@{\qquad}lr@{}}
\toprule
\textbf{Category} & \textbf{Count} & \textbf{Category} & \textbf{Count} \\
\midrule
vehicle.car                & 390 & vehicle.construction              & 6 \\
human.pedestrian.adult     & 213 & human.pedestrian.child            & 4 \\
vehicle.truck              & 28  & vehicle.trailer                   & 2 \\
vehicle.motorcycle         & 20  & human.pedestrian.police\_officer  & 1 \\
vehicle.bicycle            & 15  & human.pedestrian.personal\_mobility & 1 \\
vehicle.bus (rigid+bendy)  & 15  & human.pedestrian.construction\_worker & 8 \\
\midrule
\multicolumn{3}{l}{\textbf{Total unique agents}} & \textbf{703} \\
\bottomrule
\end{tabular}
\end{center}

\subsection{Deliverables checklist}
Every scene has completed the full chain:

\begin{center}
\small
\begin{tabular}{@{}lccccc@{}}
\toprule
\textbf{Artifact (per scene, $\times$10)} & \textbf{Status} \\
\midrule
Converted \file{.xodr} + \file{\_meta.json} (CARLA flavor)        & \checkmark \\
RoadRunner-hardened \file{.xodr} + \file{.osm} extract            & \checkmark \\
CARLA camera renders, replay and follow modes (6 cams)            & \checkmark \\
CARLA LiDAR sweeps + BEV renders                                  & \checkmark \\
Real-vs-sim comparison GIFs (cameras, LiDAR, top-down, composites)& \checkmark \\
OSM georeferencing overlay figure                                 & \checkmark \\
RoadRunner city scene (\file{\_city.rrscene} + \file{\_city.rrhd})& \checkmark \\
RoadRunner replay scenario (\file{\_replay.rrscenario})           & \checkmark \\
\bottomrule
\end{tabular}
\end{center}

\subsection{Chronology}
File timestamps trace the work: the Python conversion/validation toolchain and the
CARLA renders were completed by early July 2026 (converter and notebooks June
30--July 4); the RoadRunner track was built July 20--22 --- programmatic intersection
and direct \file{.xodr} import on July 20, OSM co-registration later that day, all ten
city scenes on July 21, and the replay scenarios finalized July 21--22 (scene-0103 and
scene-0916 re-exported on July 22).

%=====================================================================
\section{Known Limitations}
\label{sec:limitations}

\begin{itemize}
\item \textbf{Geometry.} Plan views are straight-line chains (no analytic
      arcs/spirals); curvature is only as smooth as the 0.5\,m sampling
      (Douglas--Peucker-simplified to 3\,cm tolerance for RoadRunner). Elevation is
      flat ($z=0$).
\item \textbf{Lane model.} One driving lane per road (id $-1$, laneOffset trick);
      multi-lane carriageways are represented as parallel single-lane roads.
      Singapore's left-hand traffic is emitted as right-side lanes; travel direction
      follows the nuScenes lane direction, so drivability is correct but handedness
      semantics are not modelled.
\item \textbf{Junctions.} Where connectors don't share lanes inside the crop, many
      small single-connector junctions appear --- valid OpenDRIVE, just not visually
      consolidated.
\item \textbf{Registration.} The OSM \code{manualShift} calibration has been measured
      only for scene-0103 ($[0\ 0]$); the other nine scenes use no shift and may have
      buildings a few metres off.
\item \textbf{Signals and rules.} Scripted replay/follow ego motion ignores traffic
      lights, as do the replayed agents (they follow recorded reality, which obeyed
      the real lights). Two scenes (0655, 0916) have no signals in their patches.
\item \textbf{Replay fidelity vs.\ RoadRunner constraints.} Collision pre-pruning
      deliberately truncates or drops agents whose (inflated) boxes would touch ---
      fidelity is traded for RoadRunner's fail-on-collision rule; 504 of 703 agents
      survive as full or truncated trajectories.
\item \textbf{Converter/loader details.} \code{--osm-radius} is declared but signal
      matching actually uses nearest-approach-lane; building synthesis
      (\code{--buildings} in the converter) is heuristic negative-space tiling, whereas
      the RoadRunner track uses real OSM footprints.
\end{itemize}

%=====================================================================
\section{Suggested Next Steps}
\begin{itemize}
\item Calibrate \code{manualShift} for the remaining nine scenes (the overlay figure
      workflow makes this a few minutes per scene).
\item Add elevation --- e.g.\ from LiDAR ground fits --- to the OpenDRIVE and the
      RoadRunner ground slab.
\item Consolidate small junctions and emit analytic arcs for smoother RoadRunner
      editing.
\item Close the loop between back-ends: export the RoadRunner scenarios to
      OpenSCENARIO / back into CARLA, and compare both replays frame-by-frame.
\item Use the calibrated CARLA renders for perception-model evaluation (the
      calibrated rig was designed for BEVFormer-class models) on simulated replicas of
      the real scenes, including domain-gap studies against the real camera GIFs.
\item Extend beyond v1.0-mini to the full trainval maps (the converter already supports
      any scene/version and explicit region overrides).
\end{itemize}

\appendix
%=====================================================================
\section{Appendix: Command-Line Reference}

\subsection{Converter (\file{convert\_nuscenes\_to\_xodr.py})}
\begin{center}
\footnotesize
\begin{tabular}{@{}lll@{}}
\toprule
\textbf{Flag} & \textbf{Default} & \textbf{Purpose} \\
\midrule
\code{--dataroot} / \code{--map-dataroot} & mini split path & dataset / map expansion roots \\
\code{--version}        & \code{v1.0-mini}  & nuScenes version \\
\code{--scene-index} / \code{--scene-name} & 0 / --- & scene selection \\
\code{--all-scenes}     & off               & convert every scene \\
\code{--region}         & ---               & explicit patch override (map frame) \\
\code{--padding}        & 50\,m             & margin around ego bbox \\
\code{--resolution}     & 0.5\,m            & centerline sampling step \\
\code{--lane-width}     & per-lane estimate & fixed width override \\
\code{--crosswalks}     & off               & emit crosswalk objects + polygons \\
\code{--buildings}      & off               & synthesize sidecar footprints \\
\code{--traffic-lights} & off               & emit signals + phase controllers \\
\code{--osm-signals} / \code{--osm-radius} & off / 40\,m & real OSM light counts \\
\code{--include-static} & off               & export cones/barriers/debris too \\
\code{--out}            & auto-named        & output path \\
\bottomrule
\end{tabular}
\end{center}

\subsection{CARLA loader (\file{load\_xodr\_in\_carla.py})}
\begin{center}
\scriptsize
\setlength{\tabcolsep}{3pt}
\begin{tabular}{@{}lll@{}}
\toprule
\textbf{Flag} & \textbf{Default} & \textbf{Purpose} \\
\midrule
\code{--scenario-id} / \code{--scene} / \code{--xodr} & 0 / --- / --- & scenario selection \\
\code{--list} / \code{--all-scenarios} & --- & enumerate / batch-run all \\
\code{--host} / \code{--port} & 127.0.0.1 / 2000 & server address \\
\code{--vertex-distance} / \code{--max-road-length} & 2.0 / 500 & mesh generation \\
\code{--wall-height} / \code{--extra-width} & 0.0 / 0.6 & mesh generation \\
\code{--no-ego} / \code{--no-drive} & off & load-only / spawn-only \\
\code{--replay} (default) / \code{--follow} / \code{--manual} & replay & driving mode \\
\code{--joystick}       & off    & manual: gamepad/wheel input \\
\code{--carla-agents}   & PythonAPI path & \code{agents} package for follow \\
\code{--no-agents} / \code{--max-agents} & off / 0 & agent population control \\
\code{--replay-speed} / \code{--substeps} & 1.0 / 10 & playback speed / interpolation \\
\code{--loop} (default) / \code{--no-loop} & loop & repeat vs.\ single pass \\
\code{--no-traj} / \code{--traj-hz} & on / 2\,Hz & overlay toggle / sampling \\
\code{--hist-seconds} / \code{--future-seconds} & 2 / 6 & overlay history/future span \\
\code{--cameras} / \code{--cam-dir} & off / \file{output/cameras} & camera rig \\
\code{--cam-width} / \code{--cam-height} & 800 / 450 & approx-rig resolution \\
\code{--calib-dir} / \code{--approx-rig} & phase0 path / off & calibrated vs.\ legacy rig \\
\code{--lidar} / \code{--lidar-dir} & off / \file{output/lidar\_carla} & LiDAR capture \\
\code{--lidar-range} / \code{--lidar-channels} & 70\,m / 32 & LiDAR model \\
\code{--lidar-points} & 34{,}000 & points per sweep \\
\code{--buildings} / \code{--max-buildings} / \code{--building-layers} & off / 0 / 3 & building props \\
\code{--wireframe-buildings} / \code{--no-crosswalks} & off & scenery style \\
\code{--weather}        & ClearNoon & \code{WeatherParameters} preset \\
\code{--tm-port} / \code{--speed-diff} & 8000 / 30\% & autopilot Traffic Manager \\
\bottomrule
\end{tabular}
\end{center}

\subsection{MATLAB entry points}
\begin{center}
\footnotesize
\begin{tabular}{@{}lp{9.8cm}@{}}
\toprule
\textbf{Script} & \textbf{Role} \\
\midrule
\file{insert\_intersection.m} & Proof of concept: programmatic 4-way intersection via \code{roadrunnerHDMap} + Scene Builder overlap groups. \\
\file{load\_nuscenes\_xodr.m} & Direct OpenDRIVE import of a converted scene into the project. \\
\file{load\_osm\_roads.m} & Georeferenced co-registration of nuScenes and OSM road networks; \code{manualShift} calibration workflow. \\
\file{nuscenes\_ground\_buildings.m} & Per-scene city scene: roads + real markings + assembled signals + ground slab + OSM buildings, one \file{.rrhd} import. Scene selectable as \code{scene-01..10} or real id. \\
\file{replay\_agents.m} & Per-scene replay scenario: CSV trajectories, asset mapping, collision pre-pruning, spawn/remove times; auto-builds the city scene if missing. \\
\bottomrule
\end{tabular}
\end{center}

\end{document}
